\chapter{Fine-tuning}
\label{ch:finetuning}

A model fitted across a broad dataset is a good \emph{starting point} for a
specific material family, and a poor substitute for one. Fine-tuning continues
that fit on the smaller set at a much lower learning rate: what the base model
learned about the general map is kept, and only the specialisation is learned.

\begin{lstlisting}
# 1. the base model, trained across everything you have
poraque-train --config configs/train_config.yaml

# 2. specialise it on one family
poraque-train --config configs/train_config.yaml \
    --fine-tune --pretrained models/poraque_models.pfno \
    --root data/oxides --checkpoint-dir models/oxides
\end{lstlisting}

\section{What comes from the checkpoint}

Two things are taken from the base model rather than from this run's
configuration, and both matter.

\textbf{The architecture}, inferred from the stored tensors. A \opt{width} or
\opt{modes} left over in the config that disagreed with the weights could only
load mismatched tensors, so the tensors are the authority --- the \opt{model}
section is ignored for the \emph{shape} of a fine-tuned network.

\textbf{The normalizations.} The datasets are re-pointed at the checkpoint's
transforms, discarding the ones just fitted to the new data.

\begin{pwarn}
Refitting the normalizations would rescale the network's inputs out from under
weights trained against the old scale, and the model would spend the fine-tune
relearning the scale instead of the chemistry. It also means the new data must
fall in a comparable range: fine-tuning on a family whose densities are an
order of magnitude away is transfer learning in name only.
\end{pwarn}

\section{Freezing the lifting path}

The lifting map embeds the input field into the network's channel space before
any operator acts on it. It is the most general part of the model, so it is the
part least in need of specialising --- and freezing it removes parameters a
small dataset could otherwise overfit. The cell encoder is frozen with it: that
embeds the lattice, a property of the input rather than of the mapping.

The \textbf{projection head stays trainable}. It decodes to physical units,
which is precisely what differs between material families.

\begin{lstlisting}
fine_tuning:
  freeze_lifting_layers: true
\end{lstlisting}

The run reports what was actually frozen, so it never has to be inferred:

\begin{lstlisting}
      fine-tuning from : models/poraque_models.pfno
      architecture     : inferred from the checkpoint (4,202,001 parameters)
      transforms       : taken from the checkpoint
      frozen           : lifting path, 1,392 parameters;
                         4,200,609 remain trainable
      learning rate    : 1e-05 (fine-tuning; base training uses 0.002)
\end{lstlisting}

\begin{ptip}
Frozen parameters are kept out of the optimiser entirely, not merely denied a
gradient. AdamW's decoupled weight decay applies without one, so a frozen
weight left in the optimiser would shrink towards zero every step --- quietly
undoing the pre-training the freeze was meant to preserve.
\end{ptip}

\section{Configuration}

\begin{center}
\begin{tabular}{@{}p{3.9cm}p{3.7cm}p{7.4cm}@{}}
\toprule
Key & Default & Meaning \\
\midrule
\opt{enable} & \opt{false} & Start from \opt{pretrained\_checkpoint} instead of
a fresh initialisation. \\
\opt{pretrained\_} \opt{checkpoint} & \file{models/poraque\_} \file{models.pfno}
& Base bundle to adapt. \\
\opt{learning\_rate} & \opt{1e-05} & Replaces \opt{training.learning\_rate} for
this run. \\
\opt{freeze\_lifting\_} \opt{layers} & \opt{false} & Hold the input lifting
path fixed. \\
\bottomrule
\end{tabular}
\end{center}

The fine-tuning learning rate is much smaller by design: the base rate would
walk the weights away from the solution being adapted before a small dataset
could constrain them.

\section{Output, and the \opt{.pfno} format}
\label{sec:pfno}

Poraqu\^e writes its models as \opt{.pfno} --- \emph{Poraqu\^e Fourier Neural
Operator}:

\begin{center}
\begin{tabular}{@{}lp{8.4cm}@{}}
\toprule
File & Contents \\
\midrule
\file{models/poraque\_models.pfno} & The universal model: both operators in one
file. \\
\file{models/poraque\_finetuned.pfno} & A fine-tune, specialised to one
family. \\
\bottomrule
\end{tabular}
\end{center}

The container is a \opt{torch.save} payload keyed by task, with a format tag
and metadata; the extension is a label, not a different serialisation. Nothing
inspects it when loading, so a bundle under any name loads fine.

A fine-tune is \textbf{never} written over the base model. The two coexist in
the same directory, and a run that would genuinely overwrite its own base is
refused --- before training starts, not after.

The PDF report leads its Configuration section with the fine-tuning facts and
carries a caveat at the top, so a reader cannot mistake a specialised model's
scores for a general one's. The bundle metadata records
\opt{fine\_tuned\_from}, the learning rate and whether the lifting layers were
frozen, so a checkpoint can always be traced back to its parent.

\begin{pnote}
Models written before this rename carry \opt{.pth}. If the \opt{.pfno} file is
absent and a \opt{.pth} of the same stem is present, it is used and the
substitution is announced --- an existing trained model should not become
invisible because a default filename changed. Rename it to silence the notice.
\end{pnote}
